% ===========================================================================
\section{Faces of the SVD}\label{sec:faces}
\warmth{0}\quad\whisper{One factorization, a dozen classical tools. Each is a small edit to $\Sig$.}

\begin{strand}
With $A=U\Sig V\T$ in hand, a surprising number of textbook objects are just $\Sig$, lightly
edited. Invert the nonzero singular values and you get the pseudoinverse; take their ratio and
you get the condition number; throw $\Sig$ away entirely and keep $UV\T$ and you get the nearest
orthogonal matrix. Here is the keyring.
\end{strand}

\block{The keyring: each tool is $\Sig$, edited}
\noindent\begin{tabularx}{\linewidth}{@{}l l L@{}}
\toprule
{\color{indigo}tool} & {\color{indigo}formula} & \multicolumn{1}{l}{\color{indigo}what it does}\\
\midrule
pseudoinverse $A^{+}$      & $V\Sig^{+}U\T$              & least squares: $x^\star=A^{+}b$ minimizes $\norm{Ax-b}$\\
condition number $\kappa_2(A)$ & $\fillin[1.8cm]$        & how much solving $Ax=b$ amplifies error\\
polar factor $Q$           & $UV\T$                     & the orthogonal matrix \emph{nearest} to $A$\\
PCA directions             & top columns of $V$         & principal axes; $\sigma_i^2\propto$ variance captured\\
operator / Frobenius norm  & $\sigma_1\ \ /\ \sqrt{\textstyle\sum\sigma_i^2}$ & size of $A$ as a map / as a vector\\
best rank-$k$ copy         & truncate $\Sig$ to $k$     & §\ref{sec:eckart}\\
\bottomrule
\end{tabularx}

\begin{definition}[Moore--Penrose pseudoinverse]\label{def:pinv}
$A^{+}=V\Sig^{+}U\T$, where $\Sig^{+}$ replaces each nonzero $\sigma_i$ by $\fillin[1.2cm]$ and
keeps the zeros (and transposes the shape). It is the honest inverse on the part of the space
$A$ handles cleanly: it inverts the stretch on the row space and sends the rest to $0$.
\end{definition}

\recall{When $A$ is square and invertible, $\Sig^{+}=\Sig^{-1}$ and $A^{+}=A^{-1}$. The pseudoinverse only earns its keep when $A$ is rectangular or rank-deficient.}

\begin{yourturn}
Least squares in one line. To minimize $\norm{Ax-b}_2$ over $x$, project $b$ onto the column
space and undo the stretch:
\[
  x^\star \;=\; A^{+}b \;=\; \sum_{\sigma_i>0}\ \frac{u_i\T b}{\sigma_i}\,v_i .
\]
Where does a small $\sigma_i$ make this dangerous, and which face in the keyring measures exactly
that danger? \fillin[4cm]
\recall{A tiny $\sigma_i$ blows up $u_i\T b/\sigma_i$ — the condition number $\kappa_2=\sigma_1/\sigma_r$ is the size of that hazard. Regularization (ridge / truncated SVD) tames it.}
\end{yourturn}

\block{The one idea, recovered}
\begin{strand}
Look back at the master table on the first page and check that you can now fill every row.
A single factorization, $A=U\Sig V\T$, answered all of it: the rank is the count of nonzero
$\sigma_i$, the norms and condition number are simple functions of them, the four subspaces are
the columns of $U$ and $V$ cut at the rank, the best low-rank copy keeps the top few, and the
pseudoinverse inverts what can be inverted. \keyword{Ask a matrix a question; read the answer off
its singular values.}
\end{strand}

\loose{Where to go next: the SVD is to general matrices what the spectral theorem is to symmetric ones. For \emph{huge} $A$ you never form $A\T A$; randomized SVD and Krylov methods get the top $\sigma_i$ from a few matrix-vector products. And the same idea, lifted to operators, is the start of functional analysis.}
